% !TeX root = ../../Book2.tex
\section{Bochner–Riesz 问题}
\label{b2:sec:2403}

可积核条件给出所有指数的范围，却未必是固定 \(p\) 的最小平滑量。中心指数 \(p=2\) 无需边界光滑；离开中心后，尾部与多方向集中逐渐成为障碍。

\subsection{\texorpdfstring{\(L^2\)}{L2} 理论}
\label{b2:subsec:240301}

对每个 \(\delta\ge0\)，符号的本质上确界为 \(1\)，所以
\[
 \|B_R^\delta\|_{2\to2}=1,\qquad B_R^\delta f\to f\text{ in }L^2.
\]
收敛由 Plancherel 与频率侧支配收敛直接得到。
但 \(L^2\) 范数界不控制频率边界上的点态取值，也不自动控制 \(L^p\)、\(p\ne2\)。与限制问题的联系需要额外结构。

\ProseParagraphBreak
对于一维 \(\delta=0\)，球就是区间，Hilbert 变换给出所有 \(1<p<\infty\) 上的界。高维球不再是有限个半轴投影的乘积。矩形投影可以逐坐标应用 Hilbert 变换并用 Fubini 得到有界性；球边界的法向连续变化，出现了无限方向的组合。

% 扩写定位：b2:subsec:240301
平方理论还允许量化平均逼近的速度，而且显示有无边界消失阶会改变速度结构。
\begin{proposition}\label{b2:prop:24-l2-rate}
设 \(\delta>0\)、\(0<s\leq2\)，且 \(f\in L^2\)、\((-\Delta)^{s/2}f\in L^2\)。则
\[
 \|B_R^\delta f-f\|_2
 \leq C_\delta(2\pi R)^{-s}\|(-\Delta)^{s/2}f\|_2.
\]
若 \(\delta=0\)，则对任意 \(s>0\) 同一结论成立，常数可以取一。
\end{proposition}
\begin{proof}
对 \(a\geq0\)，初等均值估计给出
\[
 |1-(1-a)_+^\delta|\leq\max\{1,\delta\}\min\{1,a\}.
\]
对 \(0<\delta<1\)，可由 \(1-(1-a)^\delta\leq a\) 的等价凹性关系验证；对 \(\delta\geq1\)，在 \(0\leq a\leq1\) 用导数上界 \(\delta\)。令 \(a=|\xi|^2/R^2\)。当 \(0<s\leq2\)，\(\min\{1,a\}\leq a^{s/2}\)，Plancherel 给出第一项。零阶硬截断的误差只在 \(|\xi|>R\)，在那里 \(\mathbf1_{\{|\xi|>R\}}\leq(|\xi|/R)^s\)，故可用任意正阶。
\end{proof}
对于 \(\delta>0\)，在一个固定紧频带内有
\(R^2[1-(1-|\xi|^2/R^2)^\delta]\to\delta|\xi|^2\)。
若非零 \(\widehat f\) 支于该频带，支配收敛于是给出
\[
 R^2(B_R^\delta f-f)\longrightarrow
 -\frac{\delta}{4\pi^2}(-\Delta)f
 \quad\text{于 }L^2.
\]
因此通常不能仅靠输入具有更多光滑性把第一项误差提高到 \(R^{-s}\)、\(s>2\)。要提高阶数，应组合不同平均以消去 Taylor 首项，或者采用在低频附近恰好等于一的另一种截断。频率边界的平滑与低频的逼近阶数是两个不同因素。

\ProseParagraphBreak

这个例子也说明“越平滑，逼近必然越快”不是无条件原则。硬截断在固定紧频带上最终误差为零；正阶平均虽然在球边界更光滑，却在球内部偏离一，因而留下低阶 Taylor 误差。硬截断的困难主要是其他 \(L^p\) 的一致有界性和逐点求和；平方逼近本身则十分直接。

\subsection{临界指标}
\label{b2:subsec:240302}

\begin{proposition}[尾部必要条件]\label{b2:prop:24-necessary}
设 \(1\le p\le\infty\)、\(\delta\ge0\)。若 \(B_1^\delta\) 在 \(L^p\) 有界，则在 \(p\ne2\) 的相关端点上必须满足
\[
 \delta>d\left|\frac1p-\frac12\right|-\frac12.
\]
当右侧为负数时，该条件不提供比 \(\delta\ge0\) 更强的限制。
\end{proposition}
\begin{proof}
选 \(\widehat f\in C_c^\infty\)，在闭单位球附近等于 \(1\)。则 \(f\) 为 Schwartz 函数，且 \(B_1^\delta f=K_\delta\)。
核渐近中的固定比例下界说明 \(K_\delta\in L^p\)、\(p<\infty\)，必须有
\(p[\delta+(d+1)/2]>d\)，即
\(\delta>d/p-(d+1)/2\)。
符号为实值，算子自伴随；若在 \(1<p<\infty\) 上有界，便也在 \(p'\) 上有界，对 \(p'\) 使用同一条件，合并得到结论。
当 \(p=1\)，直接得到 \(\delta>(d-1)/2\)。当 \(p=\infty\)，若有界，则对 \(L^1\cap L^2\) 的配对与 \(L^\infty\) 估计给出 \(L^1\) 界，故同样得到端点条件。对于 \(p=2\)，该检验给出 \(\delta>-1/2\)，所有 \(\delta\ge0\) 都满足。
\end{proof}
严格性来自径向调和级数；没有对尾部作精确下界，就不能宣称阈值等号也失败。这个检验并未排除临界函数
\(\max\{d|1/p-1/2|-1/2,0\}\) 中为零的内部指数上的硬截断；下一小节的深层几何定理才排除高维硬球。

\subsection{Bochner–Riesz 猜想}
\label{b2:subsec:240303}

经典 Bochner–Riesz 目标是：对 \(d\ge2\)、\(1<p<\infty\)，在
\[
 \delta>\max\left\{d\left|\frac1p-\frac12\right|-\frac12,\ 0\right\}
\]
时取得 \(L^p\) 有界性。这里严格的正平滑量将硬球情形排除在目标外。维数、端点形式、径向输入限制以及局部版本各是不同问题。本章以经典目标组织证明，不把未经本项目核验的当代进展范围记作定理。

\ProseParagraphBreak

下面采用\term[fefferman-qiuchengzi-dingli]{Fefferman 球乘子定理}[Fefferman ball multiplier theorem]的下述版本。

\begin{theorem}[深层结构定理：Fefferman 球乘子定理]\label{b2:thm:24-fefferman}
在 \(d\ge2\) 时，单位球特征函数为 \(L^p(\mathbb R^d)\) 乘子、
\(1<p<\infty\)，当且仅当 \(p=2\)。
\end{theorem}
\begin{proofsketch}
假设球乘子在某个非二指数有界。对球作放大、频移并取极限，配合随机符号，可推出各方向半空间投影的一致向量值界。Besicovitch 型细矩形构造使输入的平方函数质量分散，而输出在一个小测度集合中累积；对 \(p>2\) 这与向量值界矛盾。对偶性处理 \(p<2\)，切片传递处理高维。
完整的几何矩形构造及传递论证在此略去。正式原文为
\cite[Theorem 1; Section 2, pp.~330--336]{Fefferman1971Ball}。
该原文的出版扫描已核验；本段是证明路线摘要，不作为一个完整反例构造使用。
\end{proofsketch}
这个定理并不说任何正的 \(\delta\) 都足够。球乘子的失败与正平滑后的范围之间，仍隔着一个依赖维数与 \(p\) 的估计问题。

\subsection{已知范围的结构性概览}
\label{b2:subsec:240304}

\begin{theorem}[可直接证明的插值范围]\label{b2:thm:24-interpolation-range}
对 \(1<p<\infty\)，若
\[
 \delta>(d-1)\left|\frac1p-\frac12\right|,
\]
则 \(B_R^\delta\) 在 \(L^p\) 上一致有界，且随 \(R\to\infty\) 在 \(L^p\) 中逼近恒等。
\end{theorem}
\begin{proof}
取 \(\sigma>(d-1)/2\)。对符号族 \((1-|\xi|^2)_+^z\)，在
\(\operatorname{Re}z=0\) 线上有 \(L^2\) 范数 \(1\)。
在 \(\operatorname{Re}z=\sigma\) 线上，复阶 Bessel 方程与级数给出
\[
 |K_{\sigma+it}(x)|\le C e^{C(1+t^2)}
                (1+|x|)^{-\sigma-(d+1)/2}.
\]
核在大半径的界使用\chapref{b2:ch:23}复阶 ODE 能量证明（阶数实部改变，但仍固定），小半径使用 Gamma 递推和级数。右侧可积，故有 \(L^1\) 与 \(L^\infty\) 界，增长至多 Gaussian 型。
给解析族乘 \(e^{\kappa z^2}\) 消去竖线增长，应用解析插值：
在 \(\delta=\theta\sigma\) 上，
\(|1/p-1/2|=\theta/2\)。
因此任意 \(\delta>(d-1)|1/p-1/2|\) 都可选一个这样的 \(\sigma\)。
缩放给出所有半径的一致界，上一节的稠密性论证给出逼近。
\end{proof}
这是一个明确已证范围，一般窄于经典目标。它的证明只结合 \(L^2\) 与可积核端点，没有利用球面限制的额外改进。比较两条临界直线，能看清深层理论应在哪个区间增加信息。

已经证明的一致乘子界也能与输入正则性配合，给出其他 \(L^p\) 中的逼近速度。
\begin{proposition}\label{b2:prop:24-lp-rate}
设 \(1<p<\infty\)、\(\delta>0\)，且 \(B_R^\delta\) 在该 \(L^p\) 上一致有界。则对 \(0<s\leq2\)、\(R\geq1\)，
\[
 \|B_R^\delta f-f\|_p\leq C R^{-s}\|f\|_{H^{s,p}}.
\]
\end{proposition}
\begin{proof}
取 \(\chi\in C_c^\infty\)，在 \(|\xi|\leq1/4\) 上为一，支撑于 \(|\xi|<1/2\)。把输入分为 \(\chi(\xi/R)\) 的低频部与其补的高频部。记 \(g=J^sf\)。高频部的乘子
\[
 R^s(1-\chi(\xi/R))(1+4\pi^2|\xi|^2)^{-s/2}
\]
满足统一 Mihlin 界，因为它支于 \(|\xi|\geq R/4\)，在该处 \(R\leq4|\xi|\)。因此高频输入的 \(L^p\) 范数不超过 \(CR^{-s}\|g\|_p\)。对它再用 \(I-B_R^\delta\) 的一致界即可。

低频误差的归一化乘子为
\[
 R^s[1-(1-|\xi|^2/R^2)^\delta]\,
 \chi(\xi/R)(1+4\pi^2|\xi|^2)^{-s/2}.
\]
在这片频率中边界离支集有固定距离；把方括号写成
\((|\xi|^2/R^2)b_\delta(|\xi|^2/R^2)\)，其中 \(b_\delta\) 及所需导数一致有界。对 \(|\xi|\leq1\)，因 \(s\leq2\)、\(R\geq1\)，前因子 \(R^{s-2}\) 有界；对 \(1\leq|\xi|\leq R/2\)，因子可写为 \((|\xi|/R)^{2-s}\) 乘一个零阶有界光滑量。逐次链式求导便得到统一的尺度导数界。Mihlin 给出低频误差上界，合并两部分证明命题。
\end{proof}
这个论证只要求本章某条途径已经给出一致乘子界，并不把充分条件误当成必要条件。它也把两种误差分开：高频尾部由输入的光滑性控制，低频误差由平均符号在原点的 Taylor 阶数控制。若进一步构造消掉低阶 Taylor 系数的组合平均，后一限制可以改变，但需要重新核对组合的统一算子界。

\ProseParagraphBreak

例如，对固定 \(d,p\) 选择处于已证插值范围的 \(\delta\)，一个 \(W^{1,p}\) 输入至少有 \(R^{-1}\) 的范数逼近速度，一个 \(W^{2,p}\) 输入有 \(R^{-2}\) 的速度。这里只是范数误差，不能把它直接变成全参数逐点误差率；后者需要位置上的极大控制或更强的正则性与嵌入。

\ProseParagraphBreak

必要直线与插值直线的距离可以用一个数值模型说明。在 \(d=3\)、\(p=4\) 中，尾部条件要求 \(\delta>1/4\)，本章直接证明的插值范围要求 \(\delta>1/2\)。例如 \(\delta=3/4\) 在已证范围内，\(\delta=1/8\) 被必要条件排除，而 \(\delta=3/8\) 不能仅靠这两项结果判断。说它“没有被反例排除”与说它“已证明有界”有明确区别。

\ProseParagraphBreak

指数对偶把 \(p=4\) 与 \(p=4/3\) 的范围对应起来，因为球符号为实且算子在测试类上自伴随。缩放又使每个已证同指数界自动对所有半径具有同一常数。因此这里留下的空缺是平滑量与输入指数的估计空缺，而不是半径依赖尚未处理。若一篇文献只给固定空间球内的局部输出或径向输入，其命题还不能直接填入这个全空间、一般输入的空缺。

\ProseParagraphBreak

在 \(p\) 接近二时，尾部必要条件的右侧可能为负，此时它只要求原来已经约定的 \(\delta\geq0\)，看起来没有门槛。然而高维 \(\delta=0\) 已由球乘子定理排除所有非二指数。这个额外的几何障碍来自多方向叠加，并不是单个核尾部的一条更仔细积分能看出来的。正阶问题是否能在接近二的范围内降到任意小的平滑量，则要依赖更深的曲率和限制估计。

\ProseParagraphBreak

对角线上的 \(p=2\) 则没有这些困难：符号有界已经给出范数一，且所有 \(\delta\geq0\) 都随固定阶数的半径增加强收敛到恒等。这个特殊点是各种插值证明的可靠起点，但它不控制其他指数上的几何组合。把一个证明完全归结为“平方理论加插值”时，应能够说清另一个端点来自哪里，以及为什么它有足够强的范围和统一常数。
